\documentclass[11pt,a4paper]{article}

\usepackage{../shared/cathedral-preamble}

\title{%
  The Logical Foundations of the Cathedral:\\
  Type Theory, Trust, and the Metatheory\\
  of a Machine-Verified Reduction%
}
\author{
  Jason Robert Gochanour
}
\date{April 26, 2026}

\begin{document}

\maketitle

\begin{abstract}
We analyze the logical foundations of the Cathedral, a Lean~4
formalization establishing $\RH \iff d_N^2 \to 0$.
The proof inhabits the Calculus of Inductive Constructions (CIC)
extended with three kernel axioms (\lean{propext}, \lean{Classical.choice},
\lean{Quot.sound}) and four mathematical axioms on the critical path
(a PNT limit, an Abel summation bound, a Vasyunin convergence identity,
and a Hadamard--Titchmarsh zeta lower bound).
We establish the precise logical strength of the formalization:
it is provable in $\ZFC$ plus the consistency of the kernel axioms
(which is equiconsistent with $\ZFC$). We analyze the trust chain
from silicon to theorem, examine whether the proof can be made
constructive (partially: the converse but not the forward direction),
and compare the foundational requirements with formalizations in
Coq, Isabelle/HOL, and Mizar. We argue that the Cathedral's
axiomatic architecture provides a template for \emph{foundationally
transparent} formalization of open problems, where the logical
overhead is separated from the mathematical content with surgical
precision.
\end{abstract}

\tableofcontents

% ================================================================
\section{Introduction: What Lies Beneath the Proof}
% ================================================================

Every machine-verified proof rests on a stack of trust:

\begin{center}
\begin{tabular}{@{}cl@{}}
\toprule
\textbf{Layer} & \textbf{What Must Be Trusted} \\
\midrule
5 & Mathematical axioms (4 on crown path, 47 total) \\
4 & Mathlib library ($\sim$1.4M LOC) \\
3 & Lean kernel axioms (3: \lean{propext}, \lean{choice}, \lean{Quot.sound}) \\
2 & Lean type-checker ($\sim$5{,}000 LOC of C++ kernel) \\
1 & Compiler, OS, hardware \\
\bottomrule
\end{tabular}
\end{center}

The Cathedral lives at Layer~5. Its correctness depends on every
layer below. This paper examines Layers~2--4 in detail, characterizes
the logical strength of the entire stack, and asks: what exactly
are we trusting when we trust the Cathedral?

% ================================================================
\section{The Type Theory: CIC with Extensions}
\label{sec:cic}
% ================================================================

\subsection{The Core Calculus}

Lean~4 is based on a variant of the Calculus of Inductive Constructions
(CIC), a dependent type theory with:

\begin{itemize}
  \item \textbf{Dependent function types} ($\Pi$): $\forall (x : A), B(x)$.
  \item \textbf{Inductive types}: $\NN$, \lean{List}, \lean{Prop}, etc.
  \item \textbf{A universe hierarchy}: $\lean{Prop} : \lean{Type\,0} :
    \lean{Type\,1} : \cdots$
  \item \textbf{Proof irrelevance}: All terms of type $P : \lean{Prop}$
    are definitionally equal.
  \item \textbf{Large elimination}: Pattern matching on inductive types
    can produce types, not just terms.
\end{itemize}

The consistency of CIC (without classical extensions) follows from
its normalization, which was proved by Werner (1997) for the system
underlying Coq. Lean~4's variant differs in technical details
(quotient types, mutual inductives) but is believed equiconsistent.

\subsection{The Three Kernel Axioms}

Lean~4 adds three axioms to the core CIC:

\begin{definition}[\lean{propext}]
Propositional extensionality: $(P \iff Q) \to P = Q$ for
$P, Q : \lean{Prop}$.
\end{definition}

This is essential for working with sets (defined as predicates).
Without it, two logically equivalent propositions are not considered
equal, making set-theoretic reasoning extremely verbose.

\begin{definition}[\lean{Classical.choice}]
The axiom of choice (in a type-theoretic formulation): for any
nonempty type $A$, there exists an element $a : A$. Formally:
$\forall (A : \lean{Sort\,u}),\; \lean{Nonempty}\,A \to A$.
\end{definition}

This axiom implies the law of excluded middle ($P \lor \neg P$ for
all $P : \lean{Prop}$) and enables classical reasoning. It also
enables the \lean{noncomputable} keyword: definitions that use
\lean{choice} cannot be executed but can be reasoned about.

\begin{definition}[\lean{Quot.sound}]
Quotient soundness: if $a \sim b$ then $[a] = [b]$ in the quotient
type $A / \sim$. This is built into Lean's type theory as a
primitive, not derivable from other axioms.
\end{definition}

\begin{proposition}[Consistency]
\label{prop:consistency}
The system $\CIC + \lean{propext} + \lean{choice} + \lean{Quot.sound}$
is equiconsistent with $\ZFC$. That is: if $\ZFC$ is consistent,
then the extended CIC is consistent, and vice versa.
\end{proposition}

This is a folklore result in type theory. The forward direction
follows from the set-theoretic model of CIC; the reverse from
the interpretation of $\ZFC$ in CIC with large cardinals
(which are not needed here).

% ================================================================
\section{The Seven Axioms of the Cathedral}
\label{sec:axioms}
% ================================================================

The crown theorem \lean{nyman\_beurling\_equivalence} reports exactly
7 axioms via \lean{\#print axioms} (3 kernel + 4 mathematical):

\begin{center}
\begin{tabular}{@{}clcc@{}}
\toprule
\textbf{\#} & \textbf{Axiom} & \textbf{Origin} & \textbf{Tier} \\
\midrule
1 & \lean{propext} & Lean kernel & --- \\
2 & \lean{Classical.choice} & Lean kernel & --- \\
3 & \lean{Quot.sound} & Lean kernel & --- \\
4 & \lean{pnt\_mu\_log\_div\_k} & Cathedral & PNT \\
5 & \lean{covariance\_bound\_from\_mertens\_34} & Cathedral & Abel \\
6 & \lean{partial\_integral\_tends\_to\_formula} & Cathedral & Vasyunin \\
7 & \lean{rh\_zeta\_lower\_bound\_from\_zero\_counting} & Cathedral & Hadamard \\
\bottomrule
\end{tabular}
\end{center}

Axioms 1--3 are shared by every Lean~4 proof that uses Mathlib.
They are part of the logical infrastructure, not the mathematical
content. The foundationally interesting axioms are 4--7: one encodes
a PNT consequence (Tier~PNT), one an Abel summation bound
(Tier~Abel), one a Vasyunin integral convergence (Tier~Vasyunin),
and one a Hadamard--Titchmarsh zeta lower bound (Tier~Hadamard).
Notably, the converse direction uses \textbf{zero} custom axioms.

\subsection{Axiom 4: \lean{pnt\_mu\_log\_div\_k}}

This axiom asserts a consequence of the Prime Number Theorem:

\[
  \sum_{k=1}^{N} \frac{\mu(k) \log k}{k} \to -1
  \quad \text{as } N \to \infty.
\]

This is equivalent to $-(1/\zeta)'(1) = -1$ and follows from
the PNT via Abel's limit theorem.

\begin{remark}[Logical strength]
This axiom is provable in $\ZFC$ (and indeed in much weaker systems).
It requires the PNT and a forward Tauberian theorem. The axiom is a
\emph{formalization gap}: Mathlib has the PNT but lacks the forward
Wiener--Ikehara theorem needed to derive this limit.
\end{remark}

\subsection{Axiom 5: \lean{covariance\_bound\_from\_mertens\_34}}

This axiom asserts that the Mertens $x^{3/4}$ bound implies
covariance cancellation in the Gram matrix:

\[
  |M(x)| \leq C\,x^{3/4} \implies
  v^T(G_N - bb^T)v \leq \frac{C'}{\log N}
\]

for the M\"obius witness $v_k = -\mu(k)(1 - \ln k/\ln N)$.
The infrastructure (s1\_decay, s2\_decay, s3\_uniform\_bound) is proved;
the double-sum assembly remains.

\subsection{Axiom 6: \lean{partial\_integral\_tends\_to\_formula}}

This axiom asserts the convergence of piecewise integrals in the
Vasyunin matrix representation. The diagonal case is proved via FTC;
the off-diagonal case requires the Gauss digamma formula.

\subsection{Axiom 7: \lean{rh\_zeta\_lower\_bound\_from\_zero\_counting}}

This axiom provides a polynomial lower bound on $|\zeta(s)|$ for
$\re s \geq 1/2 + \varepsilon$, which is the bedrock of the Perron
contour integration chain. It follows from Hadamard's product formula
and the zero-counting function $N(T) \sim T\log T / (2\pi)$.

\subsection{The Cathedral Is Provable in ZFC}

\begin{theorem}[Foundational Location]
\label{thm:zfc}
The crown theorem $\RH \iff d_N^2 \to 0$, as stated in the
Cathedral, is provable in $\ZFC$ (indeed in much weaker systems).
The 4 mathematical axioms are formalization gaps, not foundational
gaps. The 3 kernel axioms are equiconsistent with $\ZFC$.
\end{theorem}

This means the Cathedral does not require any large cardinal axioms,
forcing arguments, or other foundational exotica. The theorem is
ordinary mathematics, verified by a machine built on ordinary
foundations.

% ================================================================
\section{Constructivity Analysis}
\label{sec:constructive}
% ================================================================

Can the Cathedral be made constructive? That is: can the proof be
carried out without \lean{Classical.choice} (and hence without the
law of excluded middle)?

\subsection{The Converse: Partially Constructive}

The converse direction ($d_N^2 \to 0 \implies \RH$) is proved via
the Rank-1 Mellin Miracle, which proceeds by contrapositive:

\begin{quote}
If $\neg\RH$, then $\exists\,\rho$ with $\zeta(\rho) = 0$ and
$\re\rho \neq 1/2$, so $d_N^2 \geq \delta > 0$ for all $N$.
\end{quote}

The contrapositive uses classical logic (specifically, the
equivalence $(\neg Q \implies \neg P) \iff (P \implies Q)$). In
constructive logic, contrapositive reasoning is valid only for
$\neg\neg$-stable propositions.

However, the \emph{bound} $d_N^2 \geq \delta > 0$ \emph{is}
constructive: given $\rho$, the bound is computed explicitly. The
non-constructive step is the existential quantifier over $\rho$ in
$\neg\RH$.

\begin{proposition}
The converse $d_N^2 \to 0 \implies \RH$ can be proved in
constructive type theory (CIC without \lean{choice}) if one
formulates $\RH$ as $\forall\,s,\;\zeta(s) = 0 \implies
\re s = 1/2$ (a universally quantified statement) rather than
as a $\neg\exists$ statement.
\end{proposition}

\subsection{The Forward: Inherently Classical}

The forward direction ($\RH \implies d_N^2 \to 0$) uses the Mertens
function bound, which is proved using contour integration and the
residue theorem. These tools are inherently classical (they rely on
the law of excluded middle for existence of convergent subsequences,
compactness arguments, etc.).

\begin{proposition}
The forward direction $\RH \implies d_N^2 \to 0$ is not known to be
provable in constructive type theory. The deepest classical ingredient
is the Montgomery--Vaughan mean value theorem, which uses the large
sieve inequality and Perron's formula.
\end{proposition}

\subsection{Use of Choice in Mathlib}

Even if the Cathedral's mathematical content could be made constructive,
Mathlib's infrastructure uses \lean{Classical.choice} pervasively:

\begin{itemize}
  \item \lean{MeasureTheory.integral} is defined using
    \lean{Classical.choice} (to select representatives).
  \item \lean{Matrix.nonsing\_inv} uses \lean{Classical.choice}
    (to decide invertibility).
  \item \lean{Real.log} and \lean{Real.sqrt} use \lean{Decidable}
    instances derived from \lean{Classical.choice}.
\end{itemize}

Removing \lean{choice} would require replacing virtually every
Mathlib API used in the Cathedral.

% ================================================================
\section{The Trust Chain}
\label{sec:trust}
% ================================================================

\subsection{The de Bruijn Criterion}

The Cathedral satisfies the \emph{de Bruijn criterion}: its
correctness depends only on a small, auditable kernel. Lean~4's
type-checker kernel is approximately 5{,}000 lines of C++ code. An
independent implementation (Lean's \emph{lean4checker} tool) can
verify the same proofs, providing redundancy.

\begin{center}
\begin{tabular}{@{}lr@{}}
\toprule
\textbf{Component} & \textbf{Size (LOC)} \\
\midrule
Lean kernel (type-checker) & $\sim$5{,}000 \\
Lean elaborator + tactics & $\sim$100{,}000 \\
Mathlib library & $\sim$1{,}400{,}000 \\
Cathedral & $\sim$36{,}738 \\
\bottomrule
\end{tabular}
\end{center}

Only the kernel (5{,}000 LOC) is in the trusted base. The elaborator,
tactics, and Mathlib could be buggy---if so, the kernel would reject
the proof. The Cathedral could be wrong---if so, \lean{lake build}
would fail.

\subsection{Potential Failure Points}

For the Cathedral to be incorrect, one of the following would need
to be true:

\begin{enumerate}
  \item \textbf{The Lean kernel has a soundness bug.} Possible but
    extremely unlikely; the kernel is small, well-tested, and has
    been independently re-implemented.
  \item \textbf{The hardware produced a wrong result.} Possible
    (cosmic rays, etc.) but mitigated by deterministic rebuilds on
    different machines.
  \item \textbf{The mathematical axioms are false.} All four axioms are
    provable in $\ZFC$ by standard methods. If any is false,
    there is a fundamental error in classical analytic number theory.
  \item \textbf{The statement was formalized incorrectly.} This is the
    most realistic risk: the Lean type \lean{RiemannHypothesis} might
    not faithfully represent the classical \RH. Mathlib's
    \lean{riemannZeta} is carefully constructed from Dirichlet series
    and analytic continuation, matching the standard definition.
\end{enumerate}

\begin{remark}[The Specification Problem]
Risk (4) cannot be fully eliminated by any formal system. The bridge
from informal mathematics to formal types is inherently informal.
However, the Cathedral's formalization of \RH{} uses Mathlib's
\lean{riemannZeta}, which has been audited by the Lean community
and matches the standard definition via the Dirichlet series
$\sum n^{-s}$ plus analytic continuation.
\end{remark}

% ================================================================
\section{Comparison with Other Proof Assistants}
\label{sec:comparison}
% ================================================================

How would the Cathedral differ if formalized in a different system?

\begin{center}
\small
\begin{tabular}{@{}lcccc@{}}
\toprule
& \textbf{Lean~4} & \textbf{Coq} & \textbf{Isabelle/HOL} & \textbf{Mizar} \\
\midrule
Type theory & CIC + ext & CIC & HOL & Set theory \\
Classical logic & Axiom & Optional & Built-in & Built-in \\
Zeta function & Mathlib & Absent & Absent & Absent \\
Matrix algebra & Mathlib & MathComp & HOL-Analysis & MML \\
Proof style & Tactic & Tactic & Isar & Declarative \\
Trust base & 5K LOC & 8K LOC & 10K LOC & 40K LOC \\
\bottomrule
\end{tabular}
\end{center}

\subsection{Coq}

A Coq formalization would use the same core type theory (CIC) but
would need to import classical axioms explicitly (via
\lean{Classical} or \lean{Require Import ClassicalEpsilon}). The
main barrier is that Coq lacks Mathlib's analysis library: there is
no \lean{riemannZeta} in Coq. The MathComp library provides excellent
algebra but limited analysis.

A Coq formalization would have one advantage: the option to identify
which parts of the proof are constructive (using Coq's extraction
mechanism). The converse direction could potentially be extracted
to a program.

\subsection{Isabelle/HOL}

Isabelle's HOL (Higher-Order Logic) is a simple type theory, not a
dependent type theory. It is classical by default, with a small
kernel and strong proof automation (\lean{auto}, \lean{sledgehammer}).
HOL-Analysis provides real analysis but lacks the complex analysis
needed for the zeta function.

An Isabelle formalization would be more automated (Sledgehammer would
close many intermediate goals) but would require significant library
development for the zeta function.

\subsection{Mizar}

Mizar uses a set-theoretic foundation (Tarski--Grothendieck set theory),
which is closest to how mathematicians actually think. The Mizar
Mathematical Library (MML) has some number theory but lacks modern
analytic tools. Mizar's declarative proof style would produce a
formalization that reads more like traditional mathematics but is
harder to develop interactively.

\subsection{Assessment}

Lean~4 is currently the only system where the Cathedral is feasible
without massive library development, because Mathlib is the only
library that provides: (a) the Riemann zeta function, (b) $L^2$
integration, (c) matrix algebra, and (d) measure theory, all in a
single coherent framework.

% ================================================================
\section{Universe Levels and Large Eliminations}
\label{sec:universes}
% ================================================================

The Cathedral uses Lean~4's universe polymorphism minimally. The
key universe considerations:

\begin{itemize}
  \item All mathematical content lives in \lean{Type 0} (the universe
    of small types). The Gram matrix indices are \lean{Fin N}, which
    is in \lean{Type 0}.
  \item \lean{Prop} is used for all logical statements. Proof
    irrelevance ensures that the specific proof term is irrelevant;
    only the type (the statement) matters.
  \item No large eliminations are used in the critical proof path.
    The proof does not construct types by pattern-matching on proof
    terms.
  \item Universe polymorphism appears only in Mathlib's generic
    infrastructure (e.g., \lean{Module R M} is universe-polymorphic).
\end{itemize}

The foundational simplicity here is deliberate: the Cathedral avoids
universe-level complications, keeping the proof in the fragment of
CIC that is best understood and most widely trusted.

% ================================================================
\section{The Axiom-Theorem Gradient}
\label{sec:gradient}
% ================================================================

The Cathedral's axiomatic architecture suggests a foundational
concept: the \emph{axiom-theorem gradient}. In a traditional
formalization, every statement is either an axiom (accepted without
proof) or a theorem (derived from axioms). The Cathedral introduces
a finer gradation:

\begin{enumerate}
  \item \textbf{Kernel axioms}: Trusted by design; part of the logic.
    (3 in the Cathedral.)
  \item \textbf{Active mathematical axioms}: Believed true, not yet
    formalized. On the critical path. (4 in the Cathedral, v10.)
  \item \textbf{Off-path axioms}: Formalization gaps that do not
    affect the crown theorem. (43 in the Cathedral.)
  \item \textbf{Theorems}: Machine-verified from axioms and Mathlib.
  \item \textbf{Archived axioms}: Previously active, now proved and
    excised. (9 graduated axioms eliminated during v3--v10.)
\end{enumerate}

This gradient makes the \emph{progress} of formalization visible.
Each axiom elimination moves a statement from category~2 or~3 to
category~4, with category~5 providing a historical record. The
gradient is not a feature of the type theory---it is a \emph{project
management} concept layered on top of the foundations.

\begin{remark}
The axiom-theorem gradient is analogous to the debt-equity spectrum
in finance: axioms are ``debt'' (promises to prove later) and theorems
are ``equity'' (fully owned mathematical knowledge). The Cathedral's
reduction from 56 crown axioms to 4 (with 9 graduations along the way)
is a deleveraging of the proof's balance sheet.
\end{remark}

% ================================================================
\section{What Would Change If RH Were Proved?}
\label{sec:if}
% ================================================================

If the Riemann Hypothesis were proved (by any means), what would
happen to the Cathedral?

\begin{enumerate}
  \item \textbf{PNT axiom becomes a theorem.} When Mathlib gains a
    forward Wiener--Ikehara Tauberian theorem,
    \lean{pnt\_mu\_log\_div\_k} closes automatically.
  \item \textbf{Covariance bound becomes a theorem.} The infrastructure
    is proved; only the double-sum assembly remains ($\sim$200 lines).
  \item \textbf{Vasyunin convergence becomes a theorem.} Diagonal is
    proved; off-diagonal needs the Gauss digamma formula in Lean.
  \item \textbf{Hadamard bound closes.} Borel--Carath\'eodory is in
    Mathlib; the gap is connecting it to the polynomial bound.
  \item \textbf{If all 4 axioms close, the crown theorem depends on
    zero mathematical axioms}---only the 3 kernel axioms. The
    biconditional $\RH \iff d_N^2 \to 0$ would join the ranks of the
    Four-Color Theorem and the Kepler Conjecture as major results
    formalized in proof assistants.
\end{enumerate}

Conversely, if \RH{} were \emph{disproved} (a zero found off the
critical line), the Cathedral's biconditional would imply
$d_N^2 \not\to 0$: the B\'aez-Duarte distance would be bounded
away from zero. This would be a constructive consequence of the
Cathedral plus the counterexample.

% ================================================================
\section{Conclusion: Foundational Transparency}
% ================================================================

The Cathedral demonstrates that large-scale formalization of frontier
mathematics can be \emph{foundationally transparent}: every assumption
is named, classified, and tracked, from the kernel axioms at the
bottom to the mathematical axioms at the top. The trust chain is
auditable. The logical strength is characterized. The constructive
content is identified.

This transparency is not automatic. It required deliberate
architectural decisions: the axiom registry, the tiered classification,
the \lean{\#print axioms} discipline, the ghost axiom audits. Without
these, the formalization would still type-check but its foundational
status would be opaque.

We propose \emph{foundational transparency} as a design principle
for future formalizations of open problems. The goal is not merely
to produce a proof that type-checks, but to produce a proof whose
\emph{foundational commitments} are visible, minimal, and precisely
documented.

\begin{quote}
\emph{A proof is not just a certificate of truth. It is a map of
trust, showing exactly where certainty ends and faith begins.}
\end{quote}

\vspace{1cm}

\begin{thebibliography}{99}

\bibitem{lean4}
L.~de~Moura et al., \emph{The Lean~4 Theorem Prover and Programming
Language}, CADE-28, 2021.

\bibitem{werner1997}
B.~Werner, \emph{Sets in Types, Types in Sets}, Proc. TACS 1997.

\bibitem{carneiro2019}
M.~Carneiro, \emph{The Type Theory of Lean}, MSc thesis,
Carnegie Mellon University, 2019.

\bibitem{debruijn1970}
N.G.~de Bruijn, \emph{The Mathematical Language AUTOMATH},
Symposium on Automatic Demonstration, Springer LNCS 125 (1970).

\bibitem{avigad2018}
J.~Avigad, \emph{The mechanization of mathematics}, Notices of the
AMS \textbf{65} (2018), 681--690.

\bibitem{buzzard2020}
K.~Buzzard, \emph{The future of mathematics?}, ICM 2022.

\bibitem{mathlib}
The Mathlib Community, \emph{Mathlib4},
\url{https://github.com/leanprover-community/mathlib4}.

\end{thebibliography}

\end{document}
